\Article{9}
Comparison of equations (7), (9) shows that the auxiliary quantities \(A,B,C\) etc. are the values of the indeterminates \(x,y,z\) etc. corresponding to these values of the indeterminates \(\xi,\eta,\zeta\) etc.: \(\xi=\mathfrak A,\ \eta=\mathfrak B,\ \zeta=\mathfrak C\) etc. It is therefore clear that we have
\[
\left.\begin{aligned}\relax
[aa]A+[ab]B+[ac]C+\text{etc.}&=\mathfrak A\\
[ab]A+[bb]B+[bc]C+\text{etc.}&=\mathfrak B\\
[ac]A+[bc]B+[cc]C+\text{etc.}&=\mathfrak C
\end{aligned}\right\}\ (12)
\]
etc. Thus, multiplying equations (10) respectively by \(\frac ap,\ \frac{a'}{p'},\ \frac{a''}{p''}\) etc.\(\cdot\) and adding, we obtain
\[
\left.\begin{aligned}
\mathfrak A&=a\varepsilon+a'\varepsilon'+a''\varepsilon''+\text{etc.}\\
\text{and in exactly the same way}\quad&\\
\mathfrak B&=b\varepsilon+b'\varepsilon'+b''\varepsilon''+\text{etc.}\\
\mathfrak C&=c\varepsilon+c'\varepsilon'+c''\varepsilon''+\text{etc.}
\end{aligned}\right\}\ (13)
\]
etc. Now since \(\mathfrak A\) is the value of the function \(X\) when the observed values are substituted for \(v,v',v''\) etc., it will readily be seen that if the corrections \(-\varepsilon,\ -\varepsilon',\ -\varepsilon''\) etc. are respectively applied to these, the function \(X\) will thereby acquire the value \(0\), and likewise the functions
\SourceLeaf{072}
\(Y,Z\) etc. will thereby be reduced to a vanishing value. In a similar way it follows from equation (11) that \(K\) is the value of the function \(u\) resulting from the same substitution.

We shall call the application of the corrections \(-\varepsilon,-\varepsilon',-\varepsilon''\) etc. to the observations the \textit{adjustment of the observations}; we have thus clearly been led to a conclusion of the greatest importance, namely, that the observations, adjusted in the way we have shown, satisfy all the conditional equations exactly, and give to any quantity depending in any way on the observations the very same value that would result from the most suitable combination of the unchanged observations. Since it is therefore impossible to recover the errors themselves \(e,e',e''\) etc. from the conditional equations, whose number is insufficient, we have at least obtained the \textit{most plausible errors}, a name by which we may designate the quantities \(\varepsilon,\varepsilon',\varepsilon''\) etc.

\Article{10}
Since the number of observations is assumed to be greater than the number of conditional equations, infinitely many other systems satisfying the conditional equations can be found besides the system of most plausible corrections \(-\varepsilon,-\varepsilon',-\varepsilon''\) etc.; it is worth investigating how these are related to that system. Thus let \(-E,-E',-E''\) etc. constitute such a system, different from the most plausible one; we shall have
\[
\begin{aligned}
aE+a'E'+a''E''+\text{etc.}&=\mathfrak A\\
bE+b'E'+b''E''+\text{etc.}&=\mathfrak B\\
cE+c'E'+c''E''+\text{etc.}&=\mathfrak C
\end{aligned}
\]
etc. Multiplying these equations respectively by \(A,B,C\) etc. and adding, we obtain with the aid of equations (10)
\[
p\varepsilon E+p'\varepsilon'E'+p''\varepsilon''E''+\text{etc.}=A\mathfrak A+B\mathfrak B+C\mathfrak C+\text{etc.}
\]
In exactly the same way, however, equations (13) give
\[
p\varepsilon\varepsilon+p'\varepsilon'\varepsilon'+p''\varepsilon''\varepsilon''+\text{etc.}=A\mathfrak A+B\mathfrak B+C\mathfrak C+\text{etc.}\quad(14)
\]
Combining these two equations, we easily deduce
\[
\begin{aligned}
&pEE+p'E'E'+p''E''E''+\text{etc.}\\
&=p\varepsilon\varepsilon+p'\varepsilon'\varepsilon'+p''\varepsilon''\varepsilon''+\text{etc.}
\\&\qquad+p(E-\varepsilon)^2+p'(E'-\varepsilon')^2+p''(E''-\varepsilon'')^2+\text{etc.}
\end{aligned}
\]
\SourceLeaf{073}
The sum \(pEE+p'E'E'+p''E''E''+\text{etc.}\) will therefore necessarily be \textit{greater} than the sum \(p\varepsilon\varepsilon+p'\varepsilon'\varepsilon'+p''\varepsilon''\varepsilon''+\text{etc.}\), which may be stated as a

\textsc{Theorem.} \textit{The sum of the squares of the corrections by which the observations can be made to agree with the conditional equations, with each square multiplied by the weight of the corresponding observation, is a minimum if the most plausible corrections are adopted.}

This is precisely the principle of least squares, from which equations (12), (10) can also easily be derived directly. Moreover, for this smallest sum, which we shall denote below by \(S\), equation (14) gives us the expression \(\mathfrak A A+B\mathfrak B+C\mathfrak C+\text{etc.}\)

\clearpage
